The solar cooker performance is predicted by two sequential models, summarised in
Figure~\ref{fig:model_flowchart}.  A Monte Carlo \textit{ray-tracing model} first computes
the absorbed flux distribution on the receiver, which is exported as a CSV file.  That
flux map is then used as the thermal input to a \textit{heat model}, which solves the
transient energy balance and produces a temperature-versus-time prediction for the cooking
vessel.

% Two-row pipeline: RT model (row 1) → Heat model (row 2)
% Apply-equations box hangs below the iteration nodes with bidirectional arrows.
% Uses only: tikz, arrows.meta, positioning, shapes.geometric, fit (all already loaded)
\begin{figure}[H]
\centering
\resizebox{\linewidth}{!}{%
\begin{tikzpicture}[
  node distance=5mm and 7mm,
  proc/.style={rectangle, draw, rounded corners=2pt,
               minimum height=8mm, minimum width=24mm,
               text centered, font=\small, fill=blue!6},
  io/.style={trapezium, trapezium left angle=75, trapezium right angle=105,
             draw, minimum height=8mm, minimum width=24mm,
             text centered, font=\small, fill=orange!10},
  arr/.style={->, >=Stealth, thick},
  every node/.append style={align=center}
]

%── Row 1: Ray-tracing model ──────────────────────────────────────────────
\node[io]                    (json)    {Scene\\JSON};
\node[proc, right=of json]   (sample)  {Sample\\primary rays};
\node[proc, right=of sample] (trace)   {Trace \&\\intersect};
\node[proc, right=of trace]  (deposit) {Deposit\\power};
\node[proc, right=of deposit](accum)   {Accumulate\\flux stats};
\node[io,   right=of accum]  (csv)     {Flux-map\\CSV};

%── Row 2: Heat model ─────────────────────────────────────────────────────
\node[io,   below=16mm of json]      (settings) {Settings \&\\parameters};
\node[proc, right=of settings]       (charging) {Charging\\iterations};
\node[proc, right=of charging]       (cooking)  {Cooking\\iterations};
\node[io,   right=of cooking]        (tvst)     {Temperature\\vs.\ time plot};

%── Apply equations box (below the two iteration nodes) ───────────────────
\node[fit=(charging)(cooking), inner sep=0pt, draw=none] (iterfit) {};
\node[proc, below=7mm of iterfit,
      minimum width=52mm, minimum height=8mm]             (apply)
  {Apply equations: middle nodes, walls, corners};

%── Row 1 arrows ─────────────────────────────────────────────────────────
\foreach \a/\b in {json/sample, sample/trace, trace/deposit,
                   deposit/accum, accum/csv}
  \draw[arr] (\a) -- (\b);

%── Connector row 1 → row 2 (csv bends down-left into charging) ───────────
\draw[arr] (csv.south) -- ++(0,-4mm) -| (charging.north);

%── Row 2 arrows ─────────────────────────────────────────────────────────
\draw[arr] (settings.east) -- (charging.west);
\draw[arr] (charging)      -- (cooking);
\draw[arr] (cooking)       -- (tvst);

%── Bidirectional arrows: iterations ↔ apply ─────────────────────────────
\draw[{Stealth}-{Stealth}, thick]
  (charging.south) -- (charging.south |- apply.north);
\draw[{Stealth}-{Stealth}, thick]
  (cooking.south)  -- (cooking.south  |- apply.north);

%── Group labels ─────────────────────────────────────────────────────────
\node[fit=(json)(sample)(trace)(deposit)(accum)(csv),
      inner sep=0pt, draw=none] (rtfit) {};
\node[above=3pt of rtfit, font=\small\itshape] {Ray-tracing model};

\node[fit=(settings)(charging)(cooking)(tvst),
      inner sep=0pt, draw=none] (hmfit) {};
\node[above=3pt of hmfit, font=\small\itshape] {Heat model};

\end{tikzpicture}%
}
\caption{Simulation pipeline: the ray-tracing model (top) exports a flux-map CSV, which
the heat model (bottom) uses to compute a temperature-vs.-time prediction.}
\label{fig:model_flowchart}
\end{figure}


\section{Ray-Tracing Model}
\label{sec:ray_tracing_model}

\subsection{Model Scope}

The simulator estimates optical performance by Monte Carlo ray tracing.  A solar source emits
many independent rays through a user-defined aperture, each ray carries a fraction of the
incident direct normal irradiance (DNI), and surfaces transform that ray by reflection,
refraction, absorption, or splitting.  The main outputs are absorbed receiver power, receiver
flux in W/m$^2$, peak concentration ratio, and encircled-power spot size.  All geometric and
optical calculations are performed in SI units with double-precision vectors.

Each ray is represented parametrically as
\begin{equation}
  \mathbf{x}(t) = \mathbf{o} + t\mathbf{d}, \qquad
  \|\mathbf{d}\| = 1, \qquad t > 0 ,
\end{equation}
where $\mathbf{o}$ is the origin, $\mathbf{d}$ is the propagation direction, and the ray also
carries optical power $P$ in watts.  For $N$ primary rays sampled over an aperture of area
$A_\mathrm{ap}$ under irradiance $G_\mathrm{DNI}$, the initial ray power is
\begin{equation}
  P_\mathrm{ray} = \frac{G_\mathrm{DNI} A_\mathrm{ap}}{N}.
  \label{eq:ray_power}
\end{equation}
Thus the sum of deposited ray powers is a Monte Carlo estimate of total absorbed power.

\subsection{Solar Source and Primary Sampling}

Primary ray origins are sampled uniformly over a disk aperture.  The default ``pillbox'' sun
model samples directions uniformly inside the solar cone of half-angle
$\theta_s = 4.65$ mrad:
\begin{equation}
  \theta = \theta_s \sqrt{\xi_1}, \qquad
  \phi = 2\pi \xi_2 ,
\end{equation}
where $\xi_1,\xi_2 \sim U(0,1)$.  The square root gives a uniform solid-angle density within
the small solar disk.  The Buie sunshape extends this by sampling from an angular radiance
profile with a disk region and a circumsolar tail.  In solid-angle form the sampled density is
proportional to
\begin{equation}
  p(\theta) \propto L(\theta)\sin\theta ,
\end{equation}
where $L(\theta)$ is the Buie radiance profile parameterized by circumsolar ratio $\chi$.

\begin{figure}[h]
  \centering
  \fbox{\parbox[c][1.15in][c]{0.88\linewidth}{\centering
  Placeholder: aperture disk, sun cone, sampled primary rays, and receiver}}
  \caption{Primary rays are sampled on an aperture and aimed according to the selected sunshape.}
\end{figure}

\subsection{Geometry and Intersections}

Surfaces are intersected in local coordinates after applying the inverse surface transform.
For a general implicit surface $F(\mathbf{x})=0$, substituting the ray gives
$F(\mathbf{o}+t\mathbf{d})=0$; the smallest valid positive root is the first hit.  The
surface normal is the normalized gradient
\begin{equation}
  \mathbf{n} = \frac{\nabla F(\mathbf{x})}{\|\nabla F(\mathbf{x})\|},
\end{equation}
oriented against the incoming ray.  For the parabolic reflector used in the dish examples,
\begin{equation}
  F(x,y,z) = x^2+y^2-4fz = 0 ,
  \qquad
  \nabla F = (2x,\,2y,\,-4f),
\end{equation}
where $f$ is focal length.  Ideal parallel rays along the optical axis reflect through the
focus $(0,0,f)$, which is also used as a validation check.

\subsection{Optical Interactions}

Specular reflection uses the standard vector law
\begin{equation}
  \mathbf{d}_r = \mathbf{d} - 2(\mathbf{d}\cdot\mathbf{n})\mathbf{n}.
  \label{eq:reflection}
\end{equation}
A perfect mirror preserves power, while a real mirror with reflectance $\rho$ returns
\begin{equation}
  P_r = \rho P_\mathrm{in}.
\end{equation}
Optional mirror slope error perturbs the surface normal by a Gaussian angular error before
applying Eq.~\eqref{eq:reflection}, modeling imperfect alignment or roughness.

Refraction follows Snell's law,
\begin{equation}
  n_1\sin\theta_i = n_2\sin\theta_t ,
\end{equation}
with vector transmitted direction
\begin{equation}
  \mathbf{d}_t =
  \eta\mathbf{d} + \left(\eta\cos\theta_i-\cos\theta_t\right)\mathbf{n},
  \qquad
  \eta=\frac{n_1}{n_2}.
\end{equation}
If $\eta^2(1-\cos^2\theta_i)\ge 1$, total internal reflection occurs and the ray is reflected.
Otherwise, unpolarized Fresnel equations split power between reflected and transmitted rays:
\begin{align}
  R_s &=
  \left(\frac{n_1\cos\theta_i-n_2\cos\theta_t}
             {n_1\cos\theta_i+n_2\cos\theta_t}\right)^2, \\
  R_p &=
  \left(\frac{n_2\cos\theta_i-n_1\cos\theta_t}
             {n_2\cos\theta_i+n_1\cos\theta_t}\right)^2, \\
  R &= \frac{1}{2}(R_s+R_p), \qquad T=1-R .
\end{align}
The dielectric model uses deterministic splitting, so the reflected child ray receives $RP$
and the transmitted child ray receives $TP$.  Absorption inside glass follows Beer--Lambert
attenuation,
\begin{equation}
  P(d) = P_0 e^{-\alpha d},
\end{equation}
where $\alpha$ is the absorption coefficient in m$^{-1}$ and $d$ is path length in the
material.  When dispersion is enabled, the refractive index can vary with wavelength using a
Sellmeier form,
\begin{equation}
  n^2(\lambda) = 1 + \sum_{k=1}^{3}\frac{B_k\lambda^2}{\lambda^2-C_k}.
\end{equation}

\subsection{Monte Carlo Estimator and Receiver Flux}

Tracing continues until the ray escapes, is absorbed, exceeds the bounce limit, or drops below
a power cutoff.  When a ray hits an absorbing receiver, its power is deposited into a grid bin.
For receiver bin $j$ with area $\Delta A_j$, the flux estimate is
\begin{equation}
  \Phi_j = \frac{1}{\Delta A_j}\sum_{i\in j} P_i
  \qquad [\mathrm{W/m^2}].
\end{equation}
The total absorbed power, peak flux, and concentration ratio are then
\begin{equation}
  P_\mathrm{abs} = \sum_j \Phi_j \Delta A_j
  = \sum_i P_i,\qquad
  \Phi_\mathrm{peak} = \max_j \Phi_j,\qquad
  C_\mathrm{peak} = \frac{\Phi_\mathrm{peak}}{G_\mathrm{DNI}} .
\end{equation}
Spot size is reported by sorting receiver bins by radius from the target center and finding
the diameter enclosing a chosen fraction of total power, for example $D_{90}$.

\begin{figure}[h]
  \centering
  \fbox{\parbox[c][1.05in][c]{0.88\linewidth}{\centering
  Placeholder: receiver heat map with peak flux, $D_{90}$ contour, and profile slices}}
  \caption{Deposited ray power is converted to a receiver flux map and concentration metrics.}
\end{figure}

The model is therefore energy based: apart from explicitly modeled mirror losses and material
absorption, ray power is conserved through reflection, refraction, and Fresnel splitting.  As
$N$ increases, random sampling error decreases with the usual Monte Carlo rate
$O(N^{-1/2})$, while deterministic physics equations define how each sampled ray transports
power through the solar cooker geometry.
